## 深度补充：TD学习算法详解

### TD(λ)的资格迹机制！

资格迹（Eligibility Traces）是TD(λ)的核心创新，它记录了每个状态-动作对在过去被访问的"痕迹"，其衰减由λ参数控制：

**数学定义**：
$$ E_t(s,a) = \gamma \lambda E_{t-1}(s,a) + \mathbf{1}(S_t=s, A_t=a) $$

**更新规则**：
$$ Q(s,a) \leftarrow Q(s,a) + \alpha \delta_t E_t(s,a) $$

其中 $\delta_t = r + \gamma Q(s',a') - Q(s,a)$ 是TD误差。

**λ参数的影响**：
- $\lambda=0$：退化为TD(0)，没有资格迹（高偏差低方差）
- $\lambda=1$：资格迹永不衰减，接近蒙特卡洛方法（低偏差高方差）  
- $\lambda \in (0,1)$：在TD(0)和蒙特卡洛之间权衡

### n-step TD的控制变种详解！

n-step TD结合了n步回报，平衡了单步TD的偏差和多步TD的方差：

**n步回报公式**：
$$ G_t^{(n)} = R_{t+1} + \gamma R_{t+2} + \cdots + \gamma^{n-1} R_{t+n} + \gamma^n Q(s_{t+n}, a_{t+n}) $$

**更新规则**：
$$ Q(s_t,a_t) \leftarrow Q(s_t,a_t) + \alpha \left[ G_t^{(n)} - Q(s_t,a_t) \right] $$

**当n=1时**就是TD(0)，**当n→∞时**接近蒙特卡洛方法。

### Expected Sarsa详解！

Expected Sarsa是Sarsa的改进版本，使用期望而不是采样下一个动作，减少方差：

**更新规则**：
$$ Q(s,a) \leftarrow Q(s,a) + \alpha \left[ r + \gamma \sum_{a'} \pi(a'|s') Q(s',a') - Q(s,a) \right] $$

**优势**：比Sarsa方差小（因为取了期望而不是采样）
**代价**：需要知道或估计策略π的分布

### Double Q-learning详解！

Double Q-learning使用两个Q网络解耦动作选择和评估，减少过估计偏差：

**更新规则**：
$$ Q_1(s,a) \leftarrow Q_1(s,a) + \alpha \left[ r + \gamma Q_2(s', \arg\max_{a'} Q_1(s',a')) - Q_1(s,a) \right] $$
（反之亦然，随机选择更新哪个网络）

**优势**：显著减少Q-learning的过估计偏差
**劣势**：需要维护两个Q表格，内存加倍

### 完整代码示例：TD(λ)实现（资格迹版）！

```python
import numpy as np

class TDLambdaAgent:
    """TD(λ)实现（资格迹版本）"""
    
    def __init__(self, n_states, n_actions, lamda=0.9, gamma=0.99, lr=0.01):
        self.Q = np.zeros((n_states, n_actions))
        self.lamda = lamda
        self.gamma = gamma
        self.lr = lr
    
    def update(self, trajectory, rewards):
        """
        trajectory: [(s0,a0), (s1,a1), ..., (s_{T-1},a_{T-1})]
        rewards: [r1, r2, ..., r_T]
        """
        T = len(trajectory)
        E = np.zeros_like(self.Q)  # 资格迹矩阵
        
        for t, ((s, a), r) in enumerate(zip(trajectory, rewards)):
            # 计算TD误差
            if t < T-1:
                s_next, a_next = trajectory[t+1]
                td_target = r + self.gamma * self.Q[s_next, a_next]
            else:
                td_target = r  # 终止状态
            
            td_error = td_target - self.Q[s, a]
            
            # 更新资格迹：E = γλE + 1(s_t=s, a_t=a)
            E *= self.gamma * self.lamda
            E[s, a] += 1.0
            
            # 更新所有状态-动作对（根据资格迹的权重）
            self.Q += self.lr * td_error * E
```

### 完整代码示例：Expected Sarsa实现！

```python
import numpy as np

class ExpectedSarsaAgent:
    """Expected Sarsa实现"""
    
    def __init__(self, n_states, n_actions, epsilon=0.1, gamma=0.99, lr=0.01):
        self.Q = np.zeros((n_states, n_actions))
        self.epsilon = epsilon
        self.gamma = gamma
        self.lr = lr
        self.n_actions = n_actions
    
    def expected_q_value(self, state):
        """计算期望Q值：∑ π(a'|s') * Q(s',a')"""
        best_action = np.argmax(self.Q[state])
        expected = 0.0
        for a in range(self.n_actions):
            if a == best_action:
                expected += (1.0 - self.epsilon + self.epsilon/self.n_actions) * self.Q[state, a]
            else:
                expected += (self.epsilon/self.n_actions) * self.Q[state, a]
        return expected
    
    def update(self, s, a, r, s_next, done):
        if done:
            td_target = r
        else:
            td_target = r + self.gamma * self.expected_q_value(s_next)
        
        td_error = td_target - self.Q[s, a]
        self.Q[s, a] += self.lr * td_error
```

### 应用场景扩展详解！

**1. 机器人导航（实时策略）**
- 机器人需要实时更新位置价值，不能等待episode结束
- TD学习可以每步更新，边走边学
- **实际案例**：扫地机器人路径学习、无人机避障、仓储机器人导航
- **为什么适合**：TD学习使用bootstrap（自举），可以单步更新，不需要等待episode结束

**2. 股票交易（高频交易）**
- 每笔交易后需要立即更新策略，不能等收盘
- TD学习使用bootstrap，可以快速适应市场变化
- **实际案例**：高频交易算法、加密货币套利、期权对冲
- **为什么适合**：金融市场变化快，需要快速反馈更新策略

**3. 游戏AI（实时策略）**
- RTS游戏中需要实时决策，不能等一局结束
- TD学习可以边玩边学，持续优化策略
- **实际案例**：星际争霸AI、文明AI、Dota 2 AI
- **为什么适合**：游戏AI需要快速决策，TD学习可以单步更新

**4. 推荐系统（实时反馈）**
- 用户每次点击/购买后需要立即更新推荐策略
- TD学习可以快速适应新的用户行为
- **实际案例**：新闻推荐、视频推荐、商品推荐
- **为什么适合**：用户行为变化快，需要快速反馈

### 更多数学推导！

**证明TD(0)等价于随机梯度下降**：

考虑损失函数：$L(V) = \mathbb{E}[(R_{t+1} + \gamma V(S_{t+1}) - V(S_t))^2$

梯度计算：
$$ \nabla_V L = -2(R_{t+1} + \gamma V(S_{t+1}) - V(S_t)) \cdot \nabla_V V(S_t) $$

由于 $V(S_t)$ 对 $V(S_t)$ 的导数是1，对 $V(S_{t+1})$ 的导数是0（半梯度，不考虑bootstrap的影响），因此：
$$ \nabla_V L = -2\delta_t \cdot \mathbf{1}(S_t) $$
（$\mathbf{1}(S_t)$ 是one-hot向量，只有状态s_t位置为1）

SGD更新：$V \leftarrow V - \frac{1}{2} \alpha \nabla_V L = V + \alpha \delta_t \cdot \mathbf{1}(S_t)$

这正是TD(0)的更新：$V(s_t) \leftarrow V(s_t) + \alpha \delta_t$

### 更多练习题！

**练习3：λ参数选择实验**
问题：设计一个实验，在某个具体任务（如CartPole）中，如何选择最优λ？

**答案与解析**：
1. 准备环境：CartPole-v1，设置相同初始条件
2. 测试多个λ值：λ ∈ {0, 0.3, 0.5, 0.7, 0. 9, 1. 0}
3. 每个λ运行1000个episode，记录平均奖励曲线
4. 绘制性能vs λ的曲线，选择平均奖励最高的λ
5. 通常：episode短的任务选大λ（接近MC），episode长的任务选小λ（接近TD(0)）

**练习4：偏差-方差权衡分析**
问题：比较TD(0)、TD(λ)、蒙特卡洛方法的偏差和方差，并解释原因。

**答案与解析**：
| 算法 | 偏差 | 方差 | 原因 |
|------|------|------|------|
| TD(0) | 高（bootstrap） | 低（单步） | 使用bootstrap，但只依赖下一个状态 |
| TD(λ) | 中（λ控制） | 中（λ控制） | 资格迹平衡了单步和多步 |
| 蒙特卡洛 | 无（真实回报） | 高（多步累加） | 回报是多个随机奖励的和 |

**练习5：n-step TD的理论分析**
问题：当n→∞时，n-step TD的偏差和方差如何变化？

**答案与解析**：
当n→∞时：
- 偏差→0：因为使用了完整的真实回报（除了最后的bootstrap项γ^n Q(s_{t+n})，当n→∞时γ^n→0）
- 方差→∞：因为回报G_t^{(n)}包含n个随机奖励，方差随n指数增长
- 因此需要在偏差和方差之间选择最优的n（通常通过交叉验证）
ENDTEX'
echo "已创建 supp_TD_full.tex，字符数: $(wc -m supp_TD_full.tex | awk '{print $1}')"